\documentclass[11pt]{article}
\usepackage[a4paper,margin=25mm]{geometry}
\usepackage{amsmath,amssymb,amsthm,mathtools,microtype,needspace}
\usepackage[colorlinks=true,linkcolor=blue,citecolor=blue,urlcolor=blue]{hyperref}
\newtheorem{theorem}{Theorem}
\newtheorem{lemma}[theorem]{Lemma}\newtheorem{lemmaS}{Lemma}[section]
\newtheorem{corollary}[theorem]{Corollary}
\newtheorem{proposition}[theorem]{Proposition}
\theoremstyle{remark}\newtheorem{remark}[theorem]{Remark}
\newcommand{\D}{\mathbb D}
\newcommand{\C}{\mathbb C}
\newcommand{\norm}[1]{\lVert #1\rVert}
\newcommand{\Rea}{\operatorname{Re}}
\newcommand{\ee}{\mathrm e}
\title{Interior-zero complexity of polynomial\\differentiation residuals}
\author{Lluis Eriksson}
\date{18 October 2026; version 2, 21 October 2026}
\begin{document}
\maketitle
\begin{abstract}
We classify the numbers of interior zeros that permit vanishing and
superalgebraic differentiation residuals for polynomials whose zeros
lie in the closed unit disk. If a degree-\(n\) polynomial has exactly
\(M\) strictly interior zeros, then, for every complex slope,
its normalized \(H^2\) residual is at least
\(n/[2\sqrt2\,16^M(n+2M+4)^{2M+2}]\).
Consequently a prescribed budget \(b_n\) permits a superalgebraically
small residual sequence if and only if \(b_n\to\infty\).
Sufficiency holds with real coefficients, normalization at the origin,
the slope fixed by that normalization, and exactly \(b_n\) interior
zeros. In contrast, one interior zero already permits residuals
\(O(n^{-1/6})\), with family-specific asymptotic constant \(\sqrt{15}/4^{4/3}\),
whereas the circle-only infimum tends to
\(1/\sqrt8\). For one interior zero we also prove a universal
\(c n^{-13/6}\) lower bound; the optimal rate remains open.
The finite-budget obstruction combines a reciprocal identity with
classical Mahler, truncated-shift and Poisson-kernel estimates, and
extends to a stable factor times a self-inversive cofactor.
For unrestricted interior counts we prove
\(\exp(-C\sqrt n)\) lower bounds for every
\(C>\sqrt{(8/3)\log17}\). The upper construction corrects both the
exact slope and degree rounding at every large degree.
The resulting refinement of exterior power-sum bounds uses a known
residual method; the negative answer to Erd\H{o}s Problem 973 was
already known. No matching fixed-budget rates or stretched-exponential
upper bounds are claimed.
\end{abstract}

\section{Statement and relation to previous work}
All logarithms are natural. The measure on the unit circle is normalized
to have mass one. Thus for \(P(z)=\sum_{j=0}^na_jz^j\),
\[
 \norm P_2^2=\frac1{2\pi}\int_0^{2\pi}|P(\ee^{it})|^2\,dt
             =\sum_{j=0}^n|a_j|^2,
 \qquad \norm P_\infty=\max_{|z|=1}|P(z)|.
\]
Write \(\mathcal S_n\) for the nonzero polynomials of exact degree \(n\)
whose zeros are in \(\overline\D\), including multiplicity and boundary
zeros. Define
\[
 s_n=\inf_{\substack{P\in\mathcal S_n\\\lambda\in\C}}
          \frac{\norm{P'-\lambda P}_2}{\norm P_2}.
\]
There is no normalization at the origin in this definition.

For an individual nonzero polynomial write
\(\rho_\lambda(P)=\norm{P'-\lambda P}_2/\norm P_2\).
Our principal classification concerns the number of strictly interior
zeros; all root counts include multiplicity.

\begin{theorem}[Finite interior-root budget]\label{thm:count}
Let \(P\in\mathcal S_n\) have exactly \(M\) zeros in \(\D\),
where \(n\ge1\) and \(0\le M\le n\). For every \(\lambda\in\C\),
\begin{equation}\label{eq:count}
 \rho_\lambda(P)\ge
 \frac{n}{2\sqrt2\,16^M(n+2M+4)^{2M+2}}.
\end{equation}
In particular, for \(n\ge2M+4\),
\(\rho_\lambda(P)\ge2^{-6M-4}n^{-2M-1}\).
Neither the exponent nor the constant is asserted to be sharp.
\end{theorem}

\begin{theorem}[Exact criterion for a growing budget]\label{thm:budget}
Let \((b_n)\) be any integer sequence with \(0\le b_n\le n\).
The following are equivalent:
\begin{enumerate}
\item \(b_n\to\infty\).
\item There exist \(P_n\in\mathcal S_n\), with at most \(b_n\)
strictly interior zeros, and complex slopes \(\lambda_n\), such that
\(\rho_{\lambda_n}(P_n)=o(n^{-A})\) for every fixed \(A>0\).
\item There exist real-coefficient \(P_n\in\mathcal S_n\), with
\(P_n(0)=1\) and exactly \(b_n\) strictly interior zeros, such that
\(\rho_{P_n'(0)}(P_n)=o(n^{-A})\) for every fixed \(A>0\).
\end{enumerate}
For residuals merely tending to zero under the ``at most'' budget,
the necessary and sufficient condition is instead
\(b_n\ge1\) eventually, with the same choice of restricted slope
available for sufficiency.
\end{theorem}

No monotonicity of the budget is assumed. The sufficiency argument
is existential: no explicit quantitative thresholds for an arbitrary
budget are supplied, and extra interior zeros may lie arbitrarily
close to the circle. Section~\ref{sec:count} proves these statements
and gives a necessary count for stretched-exponential decay.
Their ingredients include classical multiplication and restricted
Jensen estimates \cite{brunault,haagerup,defant}, and the logarithmic
derivative of a finite Blaschke product \cite[Lemma 4.3]{kovalev}.
These tools themselves are not new.

\begin{theorem}\label{thm:main}
Put \(C_0=\sqrt{(8/3)\log17}\). For every \(C>C_0\) there is \(N_C\)
such that
\begin{equation}\label{eq:asymptotic}
 s_n\ge \ee^{-C\sqrt n}\qquad(n\ge N_C).
\end{equation}
The following explicit bound also holds:
\begin{equation}\label{eq:explicit}
 s_n\ge \ee^{-3\sqrt n}\qquad(n\ge10^4).
\end{equation}
\end{theorem}

For comparison, Tan, Wang, Huang and Chen \cite[Theorem 1.1]{tan}
considered the restricted infimum
\[
 r_n=\inf_{\substack{P\in\mathcal S_n\\P(0)=1}}
       \frac{\norm{P'-P'(0)P}_2}{\norm P_2}
\]
and proved \(r_n\ge\exp(-(1+o(1))\sqrt n\log n)\).
For an individual normalized polynomial we write
\(\rho(P)=\norm{P'-P'(0)P}_2/\norm P_2\).
Theorem~\ref{thm:circle} identifies the positive limiting residual
when every root lies on the circle. Theorem~\ref{thm:one} shows
that one interior root suffices to make it vanish, establishing the
minimal zero-count threshold. These statements
concern the global weighted boundary norm, including the exact
restricted slope.

Since \(s_n\le r_n\), Theorem~\ref{thm:main} strengthens that quantitative
estimate and allows an arbitrary scalar slope.
Their maximum-point normalization, segment integration and projected
Newton identity are retained here, with attribution. The two changes
are a local analytic bound for the logarithmic derivative and a bounded
test polynomial. These remove the logarithmic factor in the exponent.

Luo, Yang and Zhu \cite[Theorem 1.1]{luo} had already proved that exterior
power sums cannot be exponentially small in the degree, answering
Erd\H{o}s Problem 973 negatively. The application in Section~\ref{sec:power}
is a quantitative refinement of that established conclusion.
No new solution of the original yes-or-no problem is asserted.

\begin{theorem}[Superalgebraic decay]\label{thm:upper}
For every \(A>0\), there are constants \(C_A,N_A\) such that
\[
 r_n\le C_A n^{-A}\qquad\text{for every integer }n\ge N_A.
\]
The polynomials can be chosen with real coefficients, exactly
\(M_A\) zeros strictly inside the unit disk, where \(M_A\) is
independent of \(n\), and all other zeros on its boundary.
The residual of these polynomials themselves, with the restricted slope
\(P_n'(0)\), is at most \(C_An^{-A}\).
The interior zeros approach 1 with distance and radial gap
\(\Theta(n^{-1/3})\). Together with
Theorem~\ref{thm:main}, this implies superalgebraic decay of both
\(s_n,r_n\), while \(s_n^{1/n},r_n^{1/n}\to1\).
\end{theorem}

The proof in Section~\ref{sec:upper} fixes the cancellation order
and finite integer multiplicities first, then lets the reciprocal root
scale grow as \(n^{1/3}\). No explicit quantitative bounds are supplied here on
\(M_A,C_A,N_A\) or uniform control as \(A\) grows.
Section~\ref{sec:radial} explains why a fixed strict root radius
precludes superalgebraic decay of a single sequence.

Approximation by logarithmic derivatives and finite moment realization
have substantial prior history. Komarov \cite{komarovrate} obtains
geometric approximation on each fixed subdisk by logarithmic
derivatives of circle-root polynomials; \cite{komarovmoments}
also realizes prescribed initial moments with unit weights.
Those interior estimates do not supply the global weighted boundary
norm used here, with a target slope growing with degree. For example,
\(1+z^n\) approximates zero logarithmically on subdisks but has
restricted residual \(n/\sqrt2\). The present contribution is the
combination of exact slope, global concentration and correction at
every large degree, rather than moment realization alone.


\section{The sharp barrier when every zero is on the circle}\label{sec:circle}
Let \(s_n^\circ\) and \(r_n^\circ\) denote the same infima as
\(s_n,r_n\), with every zero required to lie on \(|z|=1\).

\begin{theorem}\label{thm:circle}
The unrestricted and restricted circle-root residuals satisfy
\[
 \lim_{n\to\infty}s_n^\circ
 =\lim_{n\to\infty}r_n^\circ=\frac1{\sqrt8}.
\]
The lower asymptotic bound also holds in the larger class of
self-inversive polynomials. No finite-degree extremizer is asserted.
\end{theorem}

\begin{proof}
We first prove the lower bound, allowing an arbitrary complex slope.
Rotating the variable makes this slope real and nonnegative, denoted
\(\lambda\). Self-inversivity permits a constant phase choice such that
\[
 P(\ee^{it})=\ee^{int/2}p(t),\qquad p(t)\in\mathbb R,
 \qquad p(t+2\pi)=(-1)^n p(t).
\]
These operations preserve the unrestricted normalized residual.
The familiar phase derivative \(n/2\) for circle-root polynomials
appears explicitly in \cite[(1.3)--(1.4)]{erdelyi}; here we retain
the additional scalar subtraction. Direct differentiation gives
\[
 P'-\lambda P=\ee^{i(n/2-1)t}
 \bigl[(n/2-\lambda\cos t)p-i(p'+\lambda\sin t\,p)\bigr].
\]
Integration of the cross term by parts is legitimate because
\(p^2\) is periodic. Consequently the squared residual is exactly
\begin{equation}\label{eq:circle-energy}
 \frac{\int_0^{2\pi}\!
 [p'^2+(n^2/4+\lambda^2-(n+1)\lambda\cos t)p^2]\,dt}
 {\int_0^{2\pi}p^2\,dt}.
\end{equation}
Write \(E(L)\) for the lowest periodic eigenvalue of
\(-d^2/dt^2-L\cos t\), for \(L\ge0\). In odd degree, \(|p|\)
is periodic in \(H^1\) and has the same quadratic-form values as
\(p\). Thus in both parities \eqref{eq:circle-energy} is at least
\[
 g_n(\lambda)=n^2/4+\lambda^2+E((n+1)\lambda).
\]
This is a variational lower bound, not an identification of the
periodic and antiperiodic spectra.

The change of variables \(t=2x+\pi\) gives
\(E(L)=a_0(2L)/4\), in the standard Mathieu convention
\(-d^2/dx^2+2q\cos(2x)\); see \cite[\S28.2]{dlmf}.
The use of this ground state in circular variational problems is
classical; compare \cite[\S VI.A.2]{busch}. The fixed-index expansion
\cite[(28.8.1), \(m=0,s=1\)]{dlmf} is
\[
 a_0(q)=-2q+2\sqrt q-\tfrac14+O(q^{-1/2}),
 \qquad
 E(L)=-L+\sqrt{L/2}-\tfrac1{16}+O(L^{-1/2}).
\]
We must minimize uniformly over \(\lambda\), rather than applying
this expansion only at a fixed shift. The variational principle
gives \(E(L)\ge-L\) and Lipschitz continuity, so \(g_n\) is
continuous and coercive, with
\[
 g_n(\lambda)\ge
 (\lambda-(n+1)/2)^2-(2n+1)/4.
\]
The trial value \(\lambda=n/2+1/4\) gives
\(g_n=1/8+O(n^{-1})\). Every minimizer therefore has the form
\(n/2+c_n\), where \(c_n=O(\sqrt n)\). On this window
\((n+1)\lambda\asymp n^2\), and the displayed single-variable
Mathieu remainder is uniformly \(O(n^{-1})\). Expanding its square
root uniformly on the same window yields
\[
 g_n(n/2+c)=(c-1/4)^2+\tfrac18+O((c^2+1)/n).
\]
The upper bound on the minimum and absorption of the \(c_n^2/n\)
term imply \(c_n=O(1)\). Applying the expansion again proves
\(\inf_{\lambda\ge0}g_n(\lambda)=1/8+O(n^{-1})\).
This supplies the required lower limit, also for the restricted slope.

For the upper bound, write \(n=6m+k\), \(0\le k\le5\), and for
\(m\ge3\) choose
\begin{align*}
 c_{m,k}=\cos t_{m,k}&=\tfrac12+\frac{2k-1}{8m},\\
 B_t(z)&=(1+z)^4(1-2\cos t\,z+z^2),\\
 P_{m,k}(z)&=B_{t_{m,k}}(z)^m(1+z)^k.
\end{align*}
All zeros lie on the circle, \(P_{m,k}(0)=1\), and
\[
 P_{m,k}'(0)=n/2+1/4,\qquad
 P_{m,k}'(1)/P_{m,k}(1)=n/2.
\]
The limiting factor is \(B_*=(1+z)^4(1-z+z^2)\).
With \(s=\sin^2(\theta/2)\), its squared normalized modulus is
\(\cos^6(\theta/2)\cos^2(3\theta/2)\), with a unique maximum at
\(\theta=0\), and its logarithm is \(-3\theta^2+O(\theta^4)\).
More explicitly, for the perturbed factor this modulus is
\[
 H_c(s)=(1-s)^4\left(1-\frac{2s}{1-c}\right)^2.
\]
Its secondary maximum has value
\((1+c)^6/[729(1-c)^2]<1\) for \(0\le c\le3/4\).
All six parameter sequences eventually lie in this interval.
Local negative curvature and the strict outer gap are uniform as
\(c\to1/2\). They give Gaussian domination near zero and an
exponentially small outer tail.

For \(F_{m,k}=P_{m,k}'/P_{m,k}-P_{m,k}'(0)\), the identity
\[
 B_*'/B_*-3=-\frac{3z(z-1)^2}{z^3+1}
\]
shows, on putting \(\theta=u/\sqrt m\), that
\[
 F_{m,k}(\ee^{i\theta})\longrightarrow-\tfrac14+\tfrac32u^2,
 \qquad
 |P_{m,k}(\ee^{i\theta})/P_{m,k}(1)|^2\longrightarrow\ee^{-3u^2}.
\]
Indeed the constant term of \(F_{m,k}\) at 1 is exactly \(-1/4\),
its linear term is \(O(1)\), and its quadratic term is
\(-3m/2+O(1)\). A fixed zero-free neighborhood gives the local bound
\(C(1+|\theta|+m\theta^2)\); after rescaling, a polynomial times
a Gaussian dominates. On the outer interval, differentiate the
polynomial product, removing one factor. Uniformly bounded normalized
factor derivatives give an outer squared numerator
\(O(m^2q^{2m-2})\), for some \(q<1\). Thus no estimate at a
pole of the logarithmic derivative is required.
The ratio of the complete integrals tends to
\[
 \frac{\int_{\mathbb R}(-1/4+3u^2/2)^2\ee^{-3u^2}\,du}
 {\int_{\mathbb R}\ee^{-3u^2}\,du}
 =\tfrac1{16}-\tfrac34\cdot\tfrac16+\tfrac94\cdot\tfrac1{12}
 =\tfrac18.
\]
Uniformity over the six values of \(k\) covers every sufficiently
large degree. This proves the upper bound for \(r_n^\circ\), and
\(s_n^\circ\le r_n^\circ\) completes both limits.
\end{proof}


\section{A finite residual certificate}
We first give conditions that can be checked at a fixed degree.
The asymptotic and explicit choices will then be elementary.

\begin{lemma}[Normalization and separation]\label{lem:normalize}
Suppose \(P\in\mathcal S_n\), \(\lambda\in\C\), \(y>0\), and
\[
 \frac{\norm{P'-\lambda P}_2}{\norm P_2}\le\ee^{-y},
 \qquad \sqrt{n+1}\,\ee^{-y}\le1.
\]
Set
\begin{equation}\label{eq:eta}
 \gamma_n=2+\log(2\sqrt{n+1}),\qquad
 \eta=\frac{y-\gamma_n}{n}.
\end{equation}
If \(\eta>0\), there are numbers \(v_1,\ldots,v_n\) and moments
\(m_k=n^{-1}\sum_jv_j^k\) such that
\begin{equation}\label{eq:geometry}
 m_1\in\mathbb R,\quad m_1\ge\tfrac12,\qquad
 \Rea v_j\ge\tfrac12,\quad |v_j|^2\le\frac{\Rea v_j}{\eta}.
\end{equation}
They are the reciprocals \(v_j=(1-\beta_j)^{-1}\) of the complex displacements
from the zeros to a normalized maximum point.
\end{lemma}
\begin{proof}
This is the normalization and separation argument of
\cite[Sections 2.1--2.2]{tan}, with the scalar slope left arbitrary.
Choose \(|\tau|=1\) at a maximum of \(|P|\), and put
\[
 G(w)=\frac{P(\tau w)}{P(\tau)},\qquad
 \mu=\tau\lambda,\qquad E=G'-\mu G.
\]
The denominator is nonzero because \(P\) is nonzero. We have
\(G(1)=\norm G_\infty=1\), \(\norm G_2\le1\), and
\(\norm E_2\le\ee^{-y}\).
Differentiating \(|G(\ee^{it})|^2\) at its maximum shows that
\(d=G'(1)\) is real. Bernstein's polynomial inequality gives
\(|d|\le n\). The coefficient bound
\(\norm E_\infty\le\sqrt{n+1}\norm E_2\) therefore yields
\begin{equation}\label{eq:mu}
 |\mu-d|=|E(1)|\le1,\qquad |\mu|\le n+1.
\end{equation}

Let \(\beta\) be any zero of \(G\), and set \(\delta=1-\Rea\beta\).
Integrating the exact ODE along the segment from \(\beta\) to \(1\) gives
\[
 1=\int_\beta^1\ee^{\mu(1-w)}E(w)\,dw.
\]
That segment lies in \(\overline\D\). On it,
\(\Rea(\mu(1-w))\le n\delta+2\), using \(d\in\mathbb R\),
\(|d|\le n\), \(|\mu-d|\le1\), and \(|1-\beta|\le2\).
Consequently
\[
 1\le 2\sqrt{n+1}\,\ee^{n\delta+2-y},
 \qquad \delta\ge\eta.
\]
No zero equals \(1\), since \(G(1)=1\).
For \(v=(1-\beta)^{-1}\),
\[
 \Rea v=\frac{1-\Rea\beta}{|1-\beta|^2}\ge\tfrac12,\qquad
 \frac{|v|^2}{\Rea v}=\frac1{1-\Rea\beta}\le\frac1\eta.
\]
Finally \(nm_1=G'(1)/G(1)=d\), so \(m_1\) is real and at least \(1/2\).
\end{proof}

\begin{lemma}[Local analytic estimate]\label{lem:cauchy}
With the normalization above, let \(r>0\) and suppose
\begin{equation}\label{eq:q}
 q=r\sqrt{n+1}\,\exp\bigl(-y+(2n+1)r\bigr)\le\tfrac12.
\end{equation}
Then \(G\) has no zero in \(|w-1|\le r\), and, for every integer \(k\ge2\),
\begin{equation}\label{eq:moments}
 |m_k|\le
 \frac{2\sqrt{n+1}}{n}
       \exp\bigl(-y+(2n+1)r\bigr)r^{1-k}.
\end{equation}
\end{lemma}
\begin{proof}
The disk centered at \(1\) can extend beyond \(\overline\D\);
the needed control there comes from the coefficients.
For \(|u|\le r\),
\[
 |E(1+u)|\le\sqrt{n+1}(1+r)^n\norm E_2
            \le\sqrt{n+1}\,\ee^{-y+nr}.
\]
Solving the ODE from \(G(1)=1\) gives
\[
 G(1+u)=\ee^{\mu u}(1+h(u)),\qquad
 h(u)=\int_0^u\ee^{-\mu s}E(1+s)\,ds.
\]
Using the straight segment and \eqref{eq:mu}, we obtain
\(|h(u)|\le q\le1/2\).
Thus \(G\) is nonzero on that closed disk and on a neighborhood of it.
The analytic function \(F=G'/G-\mu=E/G\) satisfies
\[
 |F(1+u)|\le2\sqrt{n+1}\,\ee^{-y+(2n+1)r}\qquad(|u|\le r).
\]
Cauchy's derivative estimate and logarithmic differentiation give
\[
 \frac{|F^{(k-1)}(1)|}{(k-1)!}
 \le 2\sqrt{n+1}\,\ee^{-y+(2n+1)r}r^{1-k},\qquad
 F^{(k-1)}(1)=(-1)^{k-1}(k-1)!\,n m_k.
\]
This proves \eqref{eq:moments}. There is no restriction \(k\le n+1\):
\(F\) is an analytic rational function on the specified disk.
\end{proof}

\begin{proposition}[Finite certificate]\label{prop:certificate}
Fix \(n\ge1\), \(y,r>0\), and an integer \(K\ge2\).
Define \(\eta\) by \eqref{eq:eta} and set \(a=2\eta\).
Suppose \(0<\eta\le2\), \(\sqrt{n+1}\ee^{-y}\le1\), and
\eqref{eq:q} holds. Put
\begin{equation}\label{eq:B}
 B=\frac{2\sqrt{n+1}\,r}{Kan}
       \exp\left(-y+(2n+1)r+
                       K\log\left(1+\frac ar\right)\right).
\end{equation}
If
\begin{equation}\label{eq:certificate}
 \frac1{K\eta}+B<\frac12,
\end{equation}
then \(s_n\ge\ee^{-y}\).
\end{proposition}
\begin{proof}
Suppose an admissible \(P,\lambda\) had residual at most \(\ee^{-y}\).
For the numbers in Lemma~\ref{lem:normalize},
\[
 |1-av_j|^2
 =1-2a\Rea v_j+a^2|v_j|^2\le1.
\]
Use the bounded Ces\`aro test
\begin{equation}\label{eq:test}
 \Psi(v)=\frac{1-(1-av)^K}{Ka}
       =v+\sum_{k=2}^K
          \frac{(-1)^{k+1}}K\binom Kk a^{k-1}v^k.
\end{equation}
It has zero constant coefficient, linear coefficient one, and
\(|\Psi(v_j)|\le2/(Ka)=1/(K\eta)\).
By Lemma~\ref{lem:cauchy}, its remaining averaged terms have modulus
at most
\begin{align*}
 &\frac{2\sqrt{n+1}\,r}{Kan}\,
       \ee^{-y+(2n+1)r}
       \left[\left(1+\frac ar\right)^K-1-\frac{Ka}{r}\right]
 \le B.
\end{align*}
The triangle inequality and \(m_1\ge1/2\) now imply
\[
 \frac12\le m_1\le
       \left|\frac1n\sum_j\Psi(v_j)\right|+B
       \le\frac1{K\eta}+B,
\]
contrary to \eqref{eq:certificate}.
Each individual residual is therefore larger than \(\ee^{-y}\);
taking its infimum gives the asserted non-strict bound.
\end{proof}

\section{Choice of parameters}
\begin{proof}[Proof of \eqref{eq:asymptotic}]
Fix \(C>C_0\), and choose \(t>0\) so small that
\[
 C^2>\frac43(2+t)\log17.
\]
Set
\[
 y=C\sqrt n,\qquad r=\frac{y}{8n},\qquad
 \eta=\frac{y-\gamma_n}{n},\qquad
 K=\left\lceil\frac{2+t}{\eta}\right\rceil.
\]
For all sufficiently large \(n\), the preliminary conditions of
Proposition~\ref{prop:certificate} hold:
\(\eta\sim C/\sqrt n\), \(q\to0\), and
\(\sqrt{n+1}\ee^{-y}\to0\).
Also \(K\eta\ge2+t\), leaving a fixed positive margin in
\eqref{eq:certificate}. We have
\[
 K\sim\frac{2+t}{C}\sqrt n,\qquad \frac ar\to16,
\]
and hence
\[
 \frac{-y+(2n+1)r+K\log(1+a/r)}{\sqrt n}
 \longrightarrow -\frac{3C}{4}+\frac{2+t}{C}\log17<0.
\]
The prefactor in \eqref{eq:B} is \(O(n^{-1})\).
Thus \(B\to0\), and the finite certificate applies.
\end{proof}

\begin{proof}[Proof of \eqref{eq:explicit}]
Take \(n\ge10^4\) and choose
\[
 y=3\sqrt n,\qquad r=\frac3{8\sqrt n},\qquad
 K=\left\lceil\frac{21}{10\eta}\right\rceil.
\]
First,
\begin{equation}\label{eq:etaexplicit}
 \gamma_n\le\frac3{20}\sqrt n,\qquad
 \frac{19}{20}\frac yn\le\eta\le\frac yn.
\end{equation}
Indeed \(\gamma_n/\sqrt n\) decreases for \(n\ge10^4\), since the
sign of its derivative is that of \(n/(n+1)-\gamma_n<0\).
At \(n=10^4\), \(\gamma_n<10<15\):
\(\ee^8>2^8=256>2\sqrt{10001}\) suffices for the first inequality.
In particular \(0<\eta<2\).
The condition \(\sqrt{n+1}\ee^{-y}<1\) is immediate, and
\[
 q=\frac38\sqrt{1+\frac1n}
       \exp\left(-\frac94\sqrt n+\frac3{8\sqrt n}\right)<\frac12.
\]
Moreover \(K\eta\ge21/10\), so
\[
 \frac12-\frac1{K\eta}\ge\frac1{42}.
\]

From \eqref{eq:etaexplicit} we obtain \(a/r\le16\) and
\(K\le(14/19)\sqrt n+1\). Using \(\log17<3\),
\begin{align*}
 -y+(2n+1)r+K\log(1+a/r)
 &\le\left(-\frac94+\frac{14}{19}\log17\right)\sqrt n
        +\log17+\frac3{8\sqrt n}\\
 &<-\frac3{76}\sqrt n+3+\frac3{8\sqrt n}<0.
\end{align*}
The last expression decreases with \(n\) and is negative at \(10^4\).
For an elementary check of \(\log17<3\), the exponential series
through degree five gives \(\ee^3>18>17\).
Finally \(Ka\ge21/5\), and therefore
\[
 \frac{2\sqrt{n+1}\,r}{Kan}
 \le\frac5{28n}\sqrt{1+\frac1n}<\frac1{5n}.
\]
It follows that \(B<1/(5n)\le1/50000<1/42\).
All conditions of Proposition~\ref{prop:certificate} are satisfied.
\end{proof}

\begin{remark}
The constants \(3\) and \(10^4\) are convenient explicit choices.
The same proof with \(r=sy/n\), \(0<s<1/2\), gives the sufficient
asymptotic condition
\[
 C^2>\frac{2\log(1+2/s)}{1-2s}.
\]
We use \(s=1/8\) to state a simple constant. Neither this expression
nor \(C_0\) is asserted to be the optimal residual constant.
\end{remark}

\section{Exterior power sums and Hardy norms}\label{sec:power}
\begin{corollary}\label{cor:powers}
Let \(|z_j|\ge1\), \(1\le j\le n\), and write
\(M_n=\max_{2\le k\le n+1}|\sum_jz_j^k|\).
Then
\[
 M_n\ge\frac1n\ee^{-3\sqrt n}\qquad(n\ge10^4).
\]
For every \(C>C_0\), uniformly over all such configurations,
\(M_n\ge n^{-1}\ee^{-C\sqrt n}\) for all sufficiently large \(n\).
No distinguished point \(z_1=1\) is needed.
\end{corollary}
\begin{proof}
We use the projected Newton identity of
\cite[Proposition 3.1]{tan}. Put
\(P(t)=\prod_{j=1}^n(1-z_jt)\) and \(p_k=\sum_jz_j^k\).
Then \(P\in\mathcal S_n\), \(P(0)=1\), and \(P'(0)=-p_1\).
Logarithmic differentiation near zero, followed by coefficient
projection onto degrees at most \(n\), gives exactly
\[
 P'-P'(0)P=-\Pi_{\le n}
       \left[\left(\sum_{k=2}^{n+1}p_kt^{k-1}\right)P\right].
\]
Projection is contractive in \(H^2\), and multiplication by a
polynomial \(S\) has \(H^2\) norm at most \(\norm S_\infty\).
Thus
\[
 s_n\le\frac{\norm{P'-P'(0)P}_2}{\norm P_2}
       \le\sum_{k=2}^{n+1}|p_k|\le nM_n.
\]
Apply Theorem~\ref{thm:main}.
\end{proof}

The factor \(1/n\) is retained in these finite statements.
It can be absorbed in a leading square-root exponent with any
strictly larger constant. This corollary strengthens the scale in
\cite{tan}; the earlier qualitative conclusion of \cite{luo} remains
fully credited.

\begin{corollary}\label{cor:hp}
Uniformly for \(1\le p\le\infty\), \(P\in\mathcal S_n\), and
\(\lambda\in\C\),
\[
 \frac{\norm{P'-\lambda P}_{H^p}}{\norm P_{H^p}}
       \ge\frac{\ee^{-3\sqrt n}}{\sqrt{n+1}}\qquad(n\ge10^4).
\]
For every \(C>C_0\), the lower bound \(\ee^{-C\sqrt n}\) holds
for all sufficiently large \(n\), uniformly over the same parameters.
\end{corollary}
\begin{proof}
Write \(R=P'-\lambda P\), of degree at most \(n\).
For \(1\le p\le2\), bounding each Fourier coefficient by
\(\norm R_1\) gives
\[
 \norm R_2\le\sqrt{n+1}\norm R_1
       \le\sqrt{n+1}\norm R_p,\qquad
 \norm P_2\ge\norm P_p.
\]
For \(p\ge2\), use
\[
 \norm R_2\le\norm R_p,\qquad
 \norm P_p\le\norm P_\infty\le\sqrt{n+1}\norm P_2.
\]
In both cases \(\norm R_p/\norm P_p\ge
(\norm R_2/\norm P_2)/\sqrt{n+1}\).
For the asymptotic statement, first apply Theorem~\ref{thm:main}
with a constant strictly between \(C_0\) and \(C\), and absorb
\(\tfrac12\log(n+1)=o(\sqrt n)\).
\end{proof}

\section{Superalgebraic upper bounds}\label{sec:upper}
We prove the upper statement by cancelling a prescribed finite number of
logarithmic-derivative terms at a strict global maximum. All constants in
this section may depend on the prescribed cancellation order \(K\).
They are not asserted to be uniform as \(K\to\infty\).

\subsection{A finite positive moment cone}

\begin{lemma}[A finite positive moment cone]\label{lem:cone}
Fix \(K\ge3\). For a rational point \(\zeta=x+iy\), \(x>0\), \(y\ge0\),
define the real column
\[
 V(\zeta)=
 \begin{cases}
  (\zeta^2,\ldots,\zeta^K)^T,&y=0,\\
  (2\Rea\zeta^2,\ldots,2\Rea\zeta^K)^T,&y>0.
 \end{cases}
\]
There is a finite collection of these columns whose cone is
\(\mathbb R^{K-1}\). The collection can contain any prescribed \(K\)
distinct positive rational real points. Its column matrix \(A_\infty\)
admits a strictly positive rational vector \(\alpha^0\) satisfying
\[
 A_\infty\alpha^0=(0,3/2,0,\ldots,0)^T.
\]
\end{lemma}
\begin{proof}
\emph{(a) The dual cone is zero.}
Let \(a=(a_2,\ldots,a_K)\) be a real vector with \(a\cdot V(\zeta)\ge0\)
for all admissible rational \(\zeta\). Then \(h(z)=\Rea\sum_{j=2}^Ka_jz^j\)
is nonnegative at all such points (the factor \(2\) is irrelevant for the
sign). By continuity \(h\ge0\) on \(\{x>0,\,y\ge0\}\), and since
\(h(\bar z)=h(z)\), on all of \(\Rea z>0\). If \(a\ne0\) and
\(l\ge2\) is the largest index with \(a_l\ne0\), choose
\(\theta\in(-\pi/2,\pi/2)\) with \(a_l\cos(l\theta)<0\); this is
possible because \(l\ge2\) and \(\cos(l\theta)\) takes both signs on that
interval. Then \(h(te^{i\theta})<0\) for large \(t\), a contradiction.
Hence \(a=0\).

\emph{(b) Finitely many columns suffice.}
Let \(\mathcal C\) be the convex cone generated by all admissible columns.
By (a) the dual cone of \(\mathcal C\) is \(\{0\}\), so, by the bipolar
theorem, the closure of \(\mathcal C\) is \(\mathbb R^{K-1}\). A convex set
in \(\mathbb R^{K-1}\) has the same interior as its closure; hence
\(\mathcal C=\mathbb R^{K-1}\). In particular each vector \(\pm e_i\)
of the standard basis is a finite nonnegative combination of admissible
columns. Let \(s_1,\ldots,s_N\) be the finitely many columns used in these
\(2(K-1)\) combinations, together with the columns of the \(K\) prescribed
real points. They are rational vectors, and their cone is
\(\mathbb R^{K-1}\).

\emph{(c) A strictly positive rational kernel vector and solution.}
Write \(A=(s_1|\cdots|s_N)\). For each \(i\), the polyhedron
\(\{\gamma\in\mathbb R^N:\gamma\ge0,\ A\gamma=-s_i\}\) is nonempty by (b)
and is defined by rational data. It contains a vertex, and vertices of such
polyhedra are rational. Let \(\gamma^{(i)}\) be a rational point of it and
put \(\kappa=\sum_i(e_i+\gamma^{(i)})\). Then \(A\kappa=0\) and
every coordinate of \(\kappa\) is at least one. Moreover \(A\alpha=w\),
\(w=(0,3/2,0,\ldots,0)^T\), has a real solution because the columns span,
hence a rational solution \(\alpha^*\) (Gaussian elimination over
\(\mathbb Q\)). For a sufficiently large rational \(\Lambda\) the vector
\(\alpha^0=\alpha^*+\Lambda\kappa\) is strictly positive, rational, and
satisfies \(A\alpha^0=w\).
\end{proof}

\subsection{Estimates shared by the upper constructions}
Throughout this subsection \(P\in\mathcal S_n\) is real with \(P(0)=1\),
with zeros \(\beta_1,\ldots,\beta_n\in\overline\D\) counted with
multiplicity. Put
\[
 v_k=(1-\beta_k)^{-1},\qquad S_j=\sum_{k=1}^nv_k^j,\qquad
 G(z)=\frac{P(z)}{P(1)},\qquad w=z-1 .
\]
Since \(|1-\beta|^2\le2(1-\Rea\beta)\) for \(|\beta|\le1\), every
\(v_k\) has \(\Rea v_k\ge\tfrac12\), hence \(|v_k|\ge\tfrac12\).
Directly, \(G(1+w)=\prod_k(1+v_kw)\),
\(P'/P(1+w)=\sum_kv_k/(1+v_kw)\), and \(S_1=P'(1)/P(1)\). From
\(v/(1+vw)=\sum_{j=0}^{K-1}v^{j+1}(-w)^j+v^{K+1}(-w)^K/(1+vw)\) we get the
exact remainder formula
\begin{equation}\label{eq:remainder}
 S_2=\cdots=S_K=0\ \Longrightarrow\
 F(z):=\frac{P'(z)}{P(z)}-S_1
 =(-w)^K\sum_{k}\frac{v_k^{K+1}}{1+v_kw}
 \quad(|v_kw|<1\ \forall k).
\end{equation}

\begin{lemmaS}[Normalized moduli of the factors]\label{lem:factors}
Let \(z=\ee^{i\theta}\), \(s=\sin^2(\theta/2)\in[0,1]\), and let
\(B=(1+z)^4(1-z+z^2)\) as in the next subsection.
\begin{enumerate}
\item[(a)] \(|(1+z)/2|^2=1-s\le\ee^{-s}\); \(|(1+z)'|/|2|=\tfrac12\).
\item[(b)] \(|B(z)/B(1)|^2=(1-s)^4(1-4s)^2\le\ee^{-3s}\) and
\(|B'(z)|/|B(1)|\le3\).
\item[(c)] For real \(v>1\) and \(f_v(z)=1-z/\beta\), \(\beta=1-1/v\):
\(|f_v(z)/f_v(1)|^2=1+4v(v-1)s\le\ee^{T s}\), \(T=4v(v-1)\), and
\(|f_v'(z)|/|f_v(1)|=v\).
\item[(d)] For \(v=x+iy\), \(x>\tfrac12\), \(y>0\), \(|v|\ge2\), and
\(f_v=(1-z/\beta)(1-z/\bar\beta)\): \(|f_v(z)/f_v(1)|^2\le\ee^{Ts}\) with
\(T=4(|v|^2+|v-1|^2)\), and \(|f_v'(z)|/|f_v(1)|\le12|v|^2\).
\end{enumerate}
\end{lemmaS}
\begin{proof}
(a) is immediate. (b) Since \((1+z)(1-z+z^2)=1+z^3\), we have
\(B=(1+z)^3(1+z^3)\), \(|B'|\le3\cdot4\cdot2+3\cdot8=48\) on the circle and
\(B(1)=16\). Moreover \(|1+z|^2=4(1-s)\) and
\(|1-z+z^2|=|2\cos\theta-1|=|1-4s|\), which gives the identity.
For \(0\le s\le\frac12\) we have \(|1-4s|\le1\) and
\((1-s)^4\le\ee^{-4s}\). For \(\frac12\le s\le1\) put \(\tau=1-s\in[0,\frac12]\);
the quantity is \(\tau^4(3-4\tau)^2\) and \(\ee^{-3s}=\ee^{-3}\ee^{3\tau}\).
If \(\tau\le\frac14\), then \(\tau^4(3-4\tau)^2\le9/256<0.036<\ee^{-3}\).
If \(\frac14\le\tau\le\frac12\), then \(\tau^2(3-4\tau)\) is increasing
(its derivative is \(6\tau(1-2\tau)\ge0\)) and at most \(\frac14\), so the
quantity is at most \(\frac1{16}<0.063\), while
\(\ee^{-3}\ee^{3\tau}\ge\ee^{-9/4}>0.105\).

(c) For \(\beta=1-1/v\): \(|1-z/\beta|^2=1+\beta^{-2}-2\beta^{-1}\cos\theta\),
so the ratio to \((1-1/\beta)^2\) is
\(1+4s\beta/(1-\beta)^2=1+4v(v-1)s\), because
\(\beta/(1-\beta)^2=v(v-1)\). Also \(|f_v'|=1/\beta\) and
\(|f_v(1)|=(1-\beta)/\beta=1/(v\beta)\).

(d) Let \(r^2=|\beta|^2\), \(b=\Rea\beta\), and
\(D=|1-\beta|^2=1+r^2-2b\). Since \(f_v\) has real coefficients, on
the circle \(|f_v(z)|^2=|z+r^2/z-2b|^2/r^4\), and with
\(\cos\theta=1-2s\), \(\sin^2\theta=4s(1-s)\),
\[
 |z+r^2/z-2b|^2=\bigl((1+r^2)(1-2s)-2b\bigr)^2+4(1-r^2)^2s(1-s).
\]
Dividing by its value \(D^2\) at \(s=0\) and expanding,
\[
 |f_v(z)/f_v(1)|^2=1+\frac{4s\,[(1-r^2)^2-(1+r^2)D]}{D^2}
   +\frac{16r^2s^2}{D^2}.
\]
Now \(D=1/U\), \(r^2=V/U\), \(b=1-x/U\) with \(Y=y^2\),
\(U=|v|^2=x^2+Y\), \(V=|v-1|^2=(x-1)^2+Y\), and \(U-V=2x-1\). Thus
\begin{equation}\label{eq:pairmod}
 |f_v(z)/f_v(1)|^2=1+8(x^2-x-Y)s+16UVs^2 .
\end{equation}
(We also confirmed \eqref{eq:pairmod} numerically to 40 digits at random points.)
Since \(UV-(x^2-x-Y)^2=(2x-1)^2Y\ge0\), the right side is at most
\((1+4\sqrt{UV}\,s)^2\le\ee^{8\sqrt{UV}s}\le\ee^{4(U+V)s}\).
Finally \(|\beta|=|v-1|/|v|\ge1-1/|v|\ge\frac12\) and
\(|f_v'|\le2|\beta|^{-1}(1+|\beta|^{-1})\le12\) on the circle, while
\(|f_v(1)|=|1-1/\beta|^2=1/(|v|^2|\beta|^2)\); hence
\(|f_v'|/|f_v(1)|\le12|v|^2|\beta|^2\le12|v|^2\).
\end{proof}

\begin{lemmaS}[From local cancellation to a residual bound]\label{lem:assemble}
Let \(P\in\mathcal S_n\) be real with \(P(0)=1\) and let \(v_k,S_j,G,F\)
be as above. Let \(K\ge2\) be an integer and assume
\(S_1=P'(0)\) and \(S_2=\cdots=S_K=0\). Let \(R\ge1\),
\(\delta\in(0,\frac12]\) and \(\kappa,C_1,C_2,\Lambda>0\), and suppose
\begin{enumerate}
\item[(i)] \(|v_k|\le R/(2\delta)\) for all \(k\), and \(R\le\delta\sqrt n\);
\item[(ii)] \(\sum_k|v_k|^{K+1}\le\Lambda\) and \(\sum_k|v_k|^2\le C_1n\);
\item[(iii)] \(|G(\ee^{i\theta})|^2\le\exp(-\kappa n\sin^2(\theta/2))\)
for all \(\theta\);
\item[(iv)] \(|P'(\ee^{i\theta})|^2\le C_2n^3|P(1)|^2\exp(-\kappa n\sin^2(\theta/2))\)
for \(\delta/R\le|\theta|\le\pi\).
\end{enumerate}
Then, with \(C\) depending only on \(K,\kappa,C_1,C_2,\delta\),
\[
 \rho(P)^2=\frac{\norm{P'-P'(0)P}_2^2}{\norm P_2^2}
 \le C\Bigl(\Lambda^2n^{-K}+n^{7/2}\exp\bigl(-\kappa\delta^2n/(\pi^2R^2)\bigr)\Bigr).
\]
\end{lemmaS}
\begin{proof}
Put \(\lambda=P'(0)=S_1\). Then \(\rho(P)^2=\int|GF|^2\,d\theta
/\int|G|^2\,d\theta\), the factor \(|P(1)|^2\) cancelling, and
\(|\lambda|\le\sum|v_k|\le2\sum|v_k|^2\le2C_1n\).

\emph{Denominator.} The function \(g(\theta)=\log|G(\ee^{i\theta})|^2
=2\Rea\sum_k\log(1+v_kw)\), \(w=\ee^{i\theta}-1\), is even because \(P\)
is real, so \(g(0)=g'(0)=0\). For \(|\theta|\le\delta/R\) we have
\(|w|\le|\theta|\le\delta/R\), hence \(|v_kw|\le\frac12\) by (i). Differentiating twice
(\(|w'|=|w''|=1\)) gives
\(|g''|\le2\sum_k(2|v_k|+4|v_k|^2)\le16\sum_k|v_k|^2\le16C_1n\), using
\(|v_k|\ge\frac12\). Since \(n^{-1/2}\le\delta/R\) by (i), Taylor's formula
gives \(g(\theta)\ge-8C_1n\theta^2\ge-8C_1\) for \(|\theta|\le n^{-1/2}\), so
\(\norm G_2^2\ge\pi^{-1}\ee^{-8C_1}n^{-1/2}\).

\emph{Inner arc \(|\theta|\le\delta/R\).} By \eqref{eq:remainder} and
(i)--(ii), \(|F|\le2|w|^K\sum|v_k|^{K+1}\le2\Lambda|\theta|^K\). By (iii)
and \(\sin^2(\theta/2)\ge\theta^2/\pi^2\),
\[
 \frac1{2\pi}\int_{|\theta|\le\delta/R}|GF|^2\,d\theta
 \le\frac{2\Lambda^2}\pi\int_{\mathbb R}\theta^{2K}
 \ee^{-\kappa n\theta^2/\pi^2}\,d\theta
 =C\Lambda^2n^{-K-1/2}.
\]

\emph{Outer arc.} Here \(|GF|^2|P(1)|^2=|P'-\lambda P|^2\le
2|P'|^2+2\lambda^2|P|^2\), so by (iii)--(iv) the integrand is at most
\(Cn^3\exp(-\kappa n\sin^2(\theta/2))\), and
\(\sin^2(\theta/2)\ge\delta^2/(\pi^2R^2)\) there. This arc contributes at most
\(Cn^3\exp(-\kappa\delta^2n/(\pi^2R^2))\). Dividing both contributions by
the denominator bound gives the claim.
\end{proof}

\subsection{A fixed number of interior zeros}
For a real reciprocal variable \(v>1\), put
\(\beta=1-1/v\) and \(f_v(z)=1-z/\beta\). For \(v=x+iy\),
\(y>0\), use the conjugate-pair factor
\(f_v(z)=(1-z/\beta)(1-z/\overline\beta)\).
Write \(\nu_v=1,2\), respectively, and define
\[
 M_j(v)=\begin{cases}v^j,&v\text{ real},\\2\Rea v^j,&\Im v>0,\end{cases}
 \qquad
 d(v)=\begin{cases}-v^2/(v-1),&v\text{ real},\\
 2\Rea[-v^2/(v-1)],&\Im v>0.\end{cases}
\]
Thus \(d(v)=f_v'(0)-f_v'(1)/f_v(1)\) (we use \(f_v(0)=1\)), and
\(\Rea v>1/2\) is equivalent to \(|\beta|<1\); also
\(M_j(v)=\sum(1-\beta')^{-j}\) over the \(\nu_v\) zeros \(\beta'\) of \(f_v\).
Use the circle-root baseline
\[
 B(z)=(1+z)^4(1-z+z^2),\qquad
 \mathfrak b_j=\sum_{\beta:\,B(\beta)=0}(1-\beta)^{-j}
 =4\cdot2^{-j}+2\cos(j\pi/3).
\]
It has degree six, constant coefficient one,
\(B'(0)=B'(1)/B(1)=3\), \(\mathfrak b_1=3\), \(\mathfrak b_2=0\), and
\(\mathfrak b_3=-3/2\).

\begin{proof}[Proof of Theorem~\ref{thm:upper}]
\textbf{Data.}
Fix \(K\ge3\). Apply Lemma~\ref{lem:cone}, with \(K\) designated
distinct positive real nodes among the \(\zeta_i\). Clear the
denominators of its strictly positive weights by an integer \(Q\),
and put \(m_i=Q\alpha_i^0\). These are positive integer group
multiplicities, fixed independently of \(n\). Set
\[
 M=\sum_i\nu_i m_i,\qquad W=\sum_i m_i M_1(\zeta_i)>0.
\]
The moments of orders \(2,\ldots,K\) are now
\(\sum_im_iM_j(\zeta_i)=(0,3Q/2,0,\ldots,0)\).

\textbf{Integer parameters.}
For every large integer \(n\), define
\begin{gather*}
 R=(n/(6Q))^{1/3},\quad \epsilon=R^{-1},\quad
 c_n\equiv n-M\pmod6,\quad |c_n-2WR|\le3,\\
 \ell_n=(n-M-c_n)/6 .
\end{gather*}
Choose a nearest integer in the specified congruence class; both
\(\ell_n,c_n\) are positive eventually. Then
\(t_n=c_n/R=2W+O(\epsilon)\) and
\(u_n=\ell_n/R^3=Q(n-M-c_n)/n=Q+O(\epsilon^2)\).

\textbf{The nonlinear system and the implicit-function step.}
Let \(\xi=(\xi_i)\) vary near \(\zeta=(\zeta_i)\), varying only the \(K\)
designated real nodes; all other nodes stay at their \(\zeta_i\).
Define \(d_\epsilon(\xi)=-\xi^2/(\xi-\epsilon)\) for a real node,
and twice its real part for a pair. For \((\epsilon,t,u,\xi)\) put
\begin{align}
 \Phi_1&=t/2+\sum_i m_i d_\epsilon(\xi_i),\label{eq:defect-slope}\\
 \Phi_2&=t\epsilon/4+\sum_i m_i M_2(\xi_i),\label{eq:defect-two}\\
 \Phi_j&=\mathfrak b_j u\epsilon^{j-3}+t2^{-j}\epsilon^{j-1}
       +\sum_i m_i M_j(\xi_i),\qquad 3\le j\le K.
       \label{eq:defect-high}
\end{align}
The map \(\Phi=(\Phi_1,\ldots,\Phi_K)\) is real-analytic on a
neighborhood of \((0,2W,Q,\zeta)\), in particular across \(\epsilon=0\):
\(\epsilon\) enters only through polynomial terms and through
\(\xi-\epsilon\), which is nonzero because \(\xi\approx\zeta_i>0\).
The scaling \(R\to\infty\) has been absorbed in \(\epsilon,t,u\), so no
singular limit is involved. At \((0,2W,Q,\zeta)\) every component vanishes:
\(\Phi_1=W-W\), \(\Phi_2=\sum m_iM_2(\zeta_i)=0\),
\(\Phi_3=-\tfrac32Q+\tfrac32Q\), and \(\Phi_j=\sum m_iM_j(\zeta_i)=0\)
for \(4\le j\le K\). The Jacobian with respect to the \(K\) designated
real positions has columns
\[
 m_i(-1,2\zeta_i,3\zeta_i^2,\ldots,K\zeta_i^{K-1})^T,
\]
whose determinant is \(\pm K!\prod_im_i\prod_{i<i'}(\zeta_{i'}-\zeta_i)\ne0\)
(a scaled Vandermonde determinant). By the implicit-function theorem there
are \(\eta_0>0\) and a real-analytic map \(\Xi\) on
\(\{|(\epsilon,t-2W,u-Q)|<\eta_0\}\) with \(\Xi(0,2W,Q)=\zeta\) and
\(\Phi(\epsilon,t,u,\Xi(\epsilon,t,u))=0\); \(\Xi\) is Lipschitz on a
smaller ball. Since \((\epsilon,t_n,u_n)\to(0,2W,Q)\) and
\(|t_n-2W|+|u_n-Q|=O(\epsilon)\), for large \(n\) the point
\(\xi_{i,n}=\Xi(\epsilon,t_n,u_n)\) solves the system and
\(\xi_{i,n}=\zeta_i+O(\epsilon)\).

\textbf{The polynomials and the exact identities.}
Put \(v_{i,n}=R\xi_{i,n}\) and
\[
 P_n(z)=B(z)^{\ell_n}(1+z)^{c_n}\prod_i f_{v_{i,n}}(z)^{m_i}.
\]
For large \(n\), its degree is exactly \(6\ell_n+c_n+M=n\), its constant
coefficient is one, and its coefficients are real. Exactly \(M\)
zeros, counted with multiplicity, lie strictly inside the disk
(\(\Rea v_{i,n}=R\Rea\xi_{i,n}>\frac12\));
the other \(n-M\) lie on the circle. The interior zeros satisfy
\(|1-\beta_{i,n}|\asymp1-|\beta_{i,n}|\asymp R^{-1}\).
Now \(S_j=\sum_k(1-\beta_k)^{-j}\) is additive over factors, so
\(S_j=\mathfrak b_j\ell_n+c_n2^{-j}+\sum_im_iM_j(v_{i,n})\), and similarly the
slope difference \(P_n'(0)-P_n'(1)/P_n(1)\) equals
\(c_n/2+\sum_im_id(v_{i,n})\), because \(B\) and \((1+z)\) contribute
\(0\) and \(1/2\). Multiplying \(\Phi_1,\ldots,\Phi_K\) by
\(R,R^2,\ldots,R^K\), respectively, and using \(R^3u_n=\ell_n\),
\(Rt_n=c_n\), gives the exact identities
\[
 \lambda_n:=P_n'(0)=P_n'(1)/P_n(1)=S_1,\qquad
 S_j=0\quad(2\le j\le K).
\]
In particular integer rounding has been corrected, and
\(\lambda_n=3\ell_n+c_n/2+R\sum_i m_iM_1(\xi_{i,n})=O(n)\).

\textbf{Verification of the hypotheses of Lemma~\ref{lem:assemble}.}
Take \(\delta=\min\{\frac12,(4\max_i|\zeta_i|)^{-1}\}\); for large \(n\),
\(|v_{i,n}|\le2R\max_i|\zeta_i|\le R/(2\delta)\), while the baseline and
the factors \(1+z\) have \(|v|\le1\). Thus (i) holds (\(R\asymp n^{1/3}\)).
Next, the \(6\ell_n+c_n\le n\) circle reciprocals have \(|v|\le1\), and the
\(M\) interior ones \(|v|\le CR\); as \(n=6QR^3\le CR^{K+1}\),
\[
 \sum_k|v_k|^{K+1}\le n+CR^{K+1}\le C_KR^{K+1}=:\Lambda,
 \qquad \sum_k|v_k|^2\le n+CR^2\le C_1n,
\]
which is (ii). For (iii), Lemma~\ref{lem:factors} gives, with
\(s=\sin^2(\theta/2)\) and since for large \(n\) all
\(|v_{i,n}|\ge2\),
\[
 |G(\ee^{i\theta})|^2\le\exp\bigl[-s\,\mathcal E_n\bigr],\qquad
 \mathcal E_n=3\ell_n+c_n-\sum_im_iT(v_{i,n}) .
\]
Here \(3\ell_n+c_n=(n-M+c_n)/2\ge(n-M)/2\), and the fixed number of
interior groups costs \(\sum_im_iT(v_{i,n})=O(R^2)\). Hence
\(\mathcal E_n\ge\kappa_0n\) for some fixed \(\kappa_0>0\) and large \(n\), and
\begin{equation}\label{eq:defect-global}
 |P_n(\ee^{i\theta})/P_n(1)|^2\le\exp[-\kappa_0 n\sin^2(\theta/2)] .
\end{equation}
For (iv) differentiate the product. By the product rule,
\(P_n'\) is the sum of three kinds of terms: \(\ell_n\) times the term with one factor
\(B\) replaced by \(B'\), \(c_n\) times the same for \(1+z\), and \(m_i\) times
the same for one factor \(f_{v_i}\). Divide by \(|P_n(1)|\) and
use Lemma~\ref{lem:factors}: the replaced factor contributes
at most \(3\), \(\frac12\), \(v_i\) or \(12|v_i|^2\), respectively, while
the squared modulus of the product of all unreplaced factors
(including the other copies of the differentiated factor) is at most
\(\exp[-s(\mathcal E_n-\tau)]\) with \(\tau=3\), \(1\) or
\(-T(v_i)<0\). Hence, for large \(n\) (so that
\(\mathcal E_n-3\ge\kappa_0n/2\)),
\[
 \frac{|P_n'(\ee^{i\theta})|}{|P_n(1)|}
 \le\Bigl(3\ell_n+\frac{c_n}2+CR^2\Bigr)
 \exp\Bigl[-\frac s2(\mathcal E_n-3)\Bigr]
 \le Cn\exp\Bigl[-\frac{\kappa_0}4ns\Bigr].
\]
This is (iv) with \(\kappa=\kappa_0/2\) (and (iii) holds with
the same \(\kappa\)). No quotient at a boundary zero is used.

\textbf{Conclusion.}
Lemma~\ref{lem:assemble} gives
\[
 \rho(P_n)^2\le C_K\bigl(R^{2K+2}n^{-K}+n^{7/2}\ee^{-cn/R^2}\bigr)
 \le C_Kn^{(2-K)/3},
 \qquad
 \rho(P_n)\le C_Kn^{(2-K)/6},
\]
since \(n/R^2\asymp n^{1/3}\) and \(R^{2K+2}\asymp n^{(2K+2)/3}\).
Choose \(K\ge6A+2\). Then \(M_A=M\) is independent of \(n\),
as required, and the polynomials \(P_n\) themselves have
\(\rho(P_n)\le C_An^{-A}\). No explicit quantitative bounds are supplied here on
\(M_A\), \(C_A\), or the starting degree as functions of \(A\); the weights
produced by Lemma~\ref{lem:cone} are existential and can be large.
\end{proof}

\begin{theorem}[One interior zero suffices]\label{thm:one}
For every sufficiently large integer \(n\), there is a real
degree-\(n\) polynomial \(P_n\), with \(P_n(0)=1\), exactly one
zero strictly inside the unit disk and all remaining zeros on the
circle, such that \(\rho(P_n)=O(n^{-1/6})\).
The interior zero is real and positive, with distance and radial gap
from 1 of order \(n^{-1/3}\). In conjunction with
Theorem~\ref{thm:circle}, the minimum number of interior zeros
permitting a sequence with vanishing residual is exactly one.
No optimal rate for one interior zero is asserted.
\end{theorem}
\begin{proof}
We use two adjustable circle angles in the degree-34 baseline
\[
 B_{a,c}(z)=(1-2az+z^2)^{15}(1-2cz+z^2)^2,
 \qquad -1<a,c<1.
\]
Its slope difference and second reciprocal moment are, respectively,
\[
 D(a,c)=-30a-4c-17,\qquad
 T(a,c)=-\frac{15a}{1-a}-\frac{2c}{1-c}.
\]
Both vanish at \((a_*,c_*)=(-2/3,3/4)\). Write
\(n-1=34m+k\), \(0\le k\le33\), and let \(v=v_n>3/2\)
be the unique solution of
\begin{equation}\label{eq:single-cubic}
 v^3-\tfrac32v^2=(n-1)/4.
\end{equation}
The left side is strictly increasing on this interval, and
\(v\sim(n/4)^{1/3}\). Set \(\beta=1-1/v\in(0,1)\).
For large \(n\), \(m>0\). Define the real targets
\[
 d_n=\frac{v^2/(v-1)-k/2}{m}=O(n^{-2/3}),\qquad
 t_n=\frac{-v^2-k/4}{m}=O(n^{-1/3}).
\]
The following explicit formulas select circle parameters solving
\(D(a_n,c_n)=d_n\), \(T(a_n,c_n)=t_n\):
\begin{align*}
 L&=(51+d_n)/2,\qquad U=17-t_n,\\
 X&=\frac{UL+221+\sqrt{(UL-289)(UL-169)}}{30U},\\
 Y&=(L-15X)/2,\qquad a_n=1-X,\qquad c_n=1-Y.
\end{align*}
Indeed these equations are equivalent to
\(15X+2Y=L\), \(15/X+2/Y=U\); elimination gives
\(15UX^2-(UL+221)X+15L=0\). The positive-square-root branch
above tends to \(X=5/3,Y=1/4\), since
\(UL\to867/2\) and the square root tends to \(391/2\).
The denominators and discriminant stay positive, and
\(a_n,c_n\in(-1,1)\) eventually, with deviations
\(O(n^{-1/3})\) from \((a_*,c_*)\). These statements are
uniform over the 34 values of \(k\).

Put
\[
 P_n(z)=B_{a_n,c_n}(z)^m(1+z)^k(1-z/\beta).
\]
It has exactly the claimed degree, constant term and zero locations.
The two target equations give the exact identities
\[
 P_n'(0)=P_n'(1)/P_n(1)=:\lambda_n,
 \qquad \sum_{j=1}^n v_j^2=0,
 \qquad v_j=(1-\beta_j)^{-1}.
\]
Here the extra roots at \(-1\) contribute \(k/2\) to the slope
difference and \(k/4\) to the second moment; the interior root
contributes \(-v^2/(v-1)\) and \(v^2\), respectively.
For a conjugation-symmetric multiset of \(N\) circle roots
avoiding 1, its reciprocal moments satisfy
\begin{equation}\label{eq:circle-third}
 M_3^{\circ}=\tfrac32 M_2^{\circ}-N/4.
\end{equation}
This follows for each pair \(1/2\pm iy\), and also for the real
reciprocal \(1/2\). Conjugation symmetry is essential and holds
here. Taking \(N=n-1\) and \(M_2^{\circ}=-v^2\),
\eqref{eq:single-cubic} gives \(\sum_jv_j^3=0\) exactly.
Moreover \(\lambda_n=(n-1)/2+v=O(n)\).

We verify the hypotheses of Lemma~\ref{lem:assemble} with \(K=3\) and
\(R=v\). Here \(S_1=P_n'(0)=\lambda_n\) and \(S_2=S_3=0\) were shown above.

\emph{Global bound.} With \(s=\sin^2(\theta/2)\) and
\(h_a(s)=(1-2s/(1-a))^2\), one has
\(|(1-2az+z^2)/(2-2a)|^2=h_a(s)\) on the circle, so
\begin{equation}\label{eq:single-product}
 |P_n(\ee^{i\theta})/P_n(1)|^2
 =H_{a_n,c_n}(s)^m(1-s)^k\,[1+4v(v-1)s],
 \quad
 H_{a,c}=h_a^{15}h_c^{2},
\end{equation}
using Lemma~\ref{lem:factors}(a),(c). The limiting function is
\(H(s):=H_{a_*,c_*}(s)=(1-6s/5)^{30}(1-8s)^4\); apart from its zeros \(1/8,5/6\), its only
interior stationary point is \(5/24\), because its logarithmic derivative is
\(-36/(1-6s/5)-32/(1-8s)\). Its values there and at 1 are
\[
 H(5/24)=(3/4)^{30}(2/3)^4<10^{-4},\qquad H(1)=7^4/5^{30}<10^{-17}.
\]
We claim there are \(\eta,\kappa_1>0\) such that
\(H_{a,c}(s)\le\ee^{-\kappa_1s}\) for \(0\le s\le1\) and
\(|a-a_*|,|c-c_*|\le\eta\). Take \(\eta\le\frac1{20}\) and \(s_0=\frac1{16}\).
On \([0,s_0]\) both bases \(1-2s/(1-a)\), \(1-2s/(1-c)\) are positive
(as \(1-c\ge\frac15\), \(2s_0/(1-c)\le\frac58\)), so \(\log H_{a,c}\) is a sum of logarithms of
positive decreasing affine functions: it is concave, vanishes at \(0\), and its
derivative at \(0\) is \(-60/(1-a)-8/(1-c)\le-60\). Hence
\(H_{a,c}\le\ee^{-60s}\) there. On \([s_0,1]\), the function
\(H\) is decreasing on \([s_0,\frac18]\), has maximum
\(H(5/24)\) on \([\frac18,\frac56]\), and is increasing from \(0\) to \(H(1)\)
on \([\frac56,1]\); therefore
\(\sup_{[s_0,1]}H=H(s_0)<0.007\).
By uniform continuity in the parameters, after shrinking \(\eta\) we have
\(H_{a,c}\le\frac12\) on \([s_0,1]\), so \(H_{a,c}\le\ee^{-(\log2)s}\). This proves
the claim with \(\kappa_1=\log2\). Since \(a_n,c_n\) converge to \(a_*,c_*\),
for large \(n\), \eqref{eq:single-product} gives
\[
 |P_n(\ee^{i\theta})/P_n(1)|^2\le
 \exp[-s(\kappa_1m-4v(v-1))]\le\exp[-\kappa n\sin^2(\theta/2)],
\]
because \(m\sim n/34\) and \(v^2=O(n^{2/3})\). This is (iii) of the lemma.

\emph{Reciprocals.} A circle root \(\ee^{i\varphi}\) with \(\cos\varphi=a\)
has \(|v|^2=1/(2(1-a))\), hence \(|v|\le2\) for all roots of the two
quadratic factors (as \(1-a\ge\frac23\), \(1-c\ge\frac15\)), and \(|v|=\frac12\) at \(-1\). With
\(R=v\sim(n/4)^{1/3}\), hypothesis (i) holds for a small fixed \(\delta\),
and (ii) holds with \(\Lambda=Cn^{4/3}\), \(C_1\) fixed:
\(\sum|v_k|^4\le16n+v^4\), \(\sum|v_k|^2\le4n+v^2\).

\emph{Outer arc.} Differentiate the product
\(\Phi_a^{15m}\Phi_c^{2m}(1+z)^kf_v\), with
\(\Phi_a=1-2az+z^2\), \(\Phi_c=1-2cz+z^2\). The normalized derivatives
\(|\Phi_a'|/|\Phi_a(1)|\le2/(1-a)\le2\), likewise for \(\Phi_c\) (at most \(10\)), and
\(1/2\), \(v\) for \(1+z\) and \(f_v\). In the term from \(\Phi_a\) the remaining
product has squared modulus
\(h_a^{15m-1}h_c^{2m}\cdots\le H_{a,c}^{m-1}\,h_a^{14}h_c^{2}\cdots\),
and in the term from \(\Phi_c\) it is
\(H_{a,c}^{m-1}h_a^{15}h_c\cdots\); the extra functions \(h_a\le1\),
\(h_c\le(2/(1-c)-1)^2\le C\) are bounded. In the terms from \(1+z\) and
\(f_v\) one factor of \(H^m\) is kept whole. Using the claim above and
\((1-s)^k\le1\), \(1+4v(v-1)s\le\ee^{4v(v-1)s}\),
\[
 \frac{|P_n'(\ee^{i\theta})|^2}{|P_n(1)|^2}\le Cn^{2}v^2
 \exp[-s(\kappa_1(m-1)-4v(v-1))]\le C n^{3}\exp[-\kappa ns],
\]
after decreasing \(\kappa\) if necessary. This is (iv). No logarithmic
quotient at a boundary zero occurs.

\emph{Conclusion.} Lemma~\ref{lem:assemble} with \(K=3\),
\(\Lambda=Cn^{4/3}\), \(R=v\asymp n^{1/3}\) yields
\[
 \rho(P_n)^2\le C\bigl(n^{8/3}n^{-3}+n^{7/2}\ee^{-cn/v^2}\bigr)
 =O(n^{-1/3}),
\]
as required. (The local remainder is \(F_n(z)=-(z-1)^3\sum_jv_j^4/(1+v_j(z-1))\) by
\eqref{eq:remainder}.)

Finally, Theorem~\ref{thm:circle} excludes vanishing residuals with
zero interior roots; this construction permits them with one.
This proves the zero-count threshold, not a sharp rate or the
superalgebraic conclusion with one root. No explicit starting
degree is claimed. The next proposition identifies the asymptotic
norm constant of this particular family.
\end{proof}


\begin{proposition}[Asymptotic constant for the explicit family]\label{prop:prefactor}
For the polynomials constructed in Theorem~\ref{thm:one},
through all sufficiently large integer degrees,
\[
 \rho(P_n)\sim\frac{\sqrt{15}}{4^{4/3}}n^{-1/6}.
\]
This constant describes that family, not a universal one-root optimum.
\end{proposition}
\begin{proof}
Retain the construction's notation and put
\(G_n=P_n/P_n(1)\), \(F_n=P_n'/P_n-\lambda_n\), and \(w=z-1\).
Its reciprocal zeros satisfy \(S_1=\lambda_n=(n-1)/2+v\),
\(S_2=S_3=0\), and \(v/n^{1/3}\to4^{-1/3}\), by \eqref{eq:single-cubic};
all boundary reciprocals have modulus at most 2. By
\eqref{eq:remainder} with \(K=3\), whenever \(|v_jw|<1\),
\[
 F_n(z)=-w^3\sum_j\frac{v_j^4}{1+v_jw}.
\]
Put \(z=\exp(iu/\sqrt n)\) with \(u\) fixed. Then \(w=iu/\sqrt n+O(1/n)\) and
\(vw=O(n^{-1/6})\to0\), so the boundary contribution to the sum is \(O(n)\)
and the single interior one is \(v^4(1+o(1))\sim4^{-4/3}n^{4/3}\). Hence
\begin{equation}\label{eq:prefactor-profile}
 n^{1/6}F_n(\ee^{iu/\sqrt n})\longrightarrow
       i4^{-4/3}u^3 .
\end{equation}
Next, \(\log G_n(1+w)=\sum_j\log(1+v_jw)\), and
\(|\log(1+x)-x+x^2/2-x^3/3|\le|x|^4/2\) for \(|x|\le\frac12\). Using
\(S_2=S_3=0\),
\[
 \log G_n(1+w)=\lambda_nw+r_n,\qquad
 |r_n|\le\tfrac12\sum_j|v_j|^4|w|^4=O(n^{4/3}n^{-2})\to0 .
\]
Since \(\Rea(\lambda_nw)=-2\lambda_n\sin^2(\theta/2)\) with
\(\theta=u/\sqrt n\), \(\sin^2(\theta/2)=u^2/(4n)(1+O(n^{-1}))\), and
\(\lambda_n/n\to1/2\), we get
\begin{equation}\label{eq:prefactor-weight}
 |G_n(\ee^{iu/\sqrt n})|^2\longrightarrow\ee^{-u^2/2}.
\end{equation}

For the domination we use the verification of
Lemma~\ref{lem:assemble} in the proof of Theorem~\ref{thm:one}.
With \(\theta=u/\sqrt n\), (iii) gives
\(|G_n|^2\le\exp(-\kappa u^2/\pi^2)\) for all \(|u|\le\pi\sqrt n\); on the inner arc
\(|\theta|\le\delta/v\), that is \(|u|\le\delta\sqrt n/v\asymp n^{1/6}\), the
bound \(|F_n|\le2\Lambda|\theta|^3\) with \(\Lambda=Cn^{4/3}\) gives
\(n^{1/3}|F_n|^2\le C u^6\). So the functions
\(|G_n|^2\) and \(n^{1/3}|G_nF_n|^2\), extended by zero beyond the inner arc,
are dominated by \(\ee^{-\kappa u^2/\pi^2}\) and \(Cu^6\ee^{-\kappa u^2/\pi^2}\), and the
inner arc exhausts \(\mathbb R\) as \(n\to\infty\). The remaining outer arc
contributes to \(n^{5/6}\norm{G_nF_n}_2^2\) at most
\(n^{5/6}\cdot Cn^{3}\exp(-cn/v^2)\to0\), by (iv) and the outer-arc estimate of
Lemma~\ref{lem:assemble}. Since
\(\norm{G_n}_2^2=(2\pi\sqrt n)^{-1}\int|G_n|^2\,du\) and
\(n^{5/6}\norm{G_nF_n}_2^2=(2\pi)^{-1}\int n^{1/3}|G_nF_n|^2\,du\) (inner arc),
dominated convergence with \eqref{eq:prefactor-profile}--\eqref{eq:prefactor-weight} yields
\[
 \sqrt n\norm{G_n}_2^2\longrightarrow\frac1{\sqrt{2\pi}},
 \qquad
 n^{5/6}\norm{G_nF_n}_2^2\longrightarrow
       \frac{15\,4^{-8/3}}{\sqrt{2\pi}},
\]
where 15 is the sixth standard Gaussian moment
(\(\int u^6\ee^{-u^2/2}du=15\sqrt{2\pi}\)).
Dividing gives \(n^{1/3}\rho(P_n)^2\to15\cdot4^{-8/3}\), i.e. the claim.
All constants in the preceding construction
are uniform over the finite set of 34 degree remainders, and the limit does
not depend on \(k\), so the limit holds through all sufficiently large
integers.
\end{proof}

\section{Interior-root budgets and their obstruction}\label{sec:count}
We first prove the all-degree bound and then the existence criterion.
The proof also identifies the structural role of the circle roots.

\subsection{Multiplication by a polynomial of small degree}
For \(T=c\prod_j(z-\gamma_j)\), write
\(\mathcal M(T)=|c|\prod_j\max(1,|\gamma_j|)\) for its Mahler measure.
The classical coefficient bound gives
\begin{equation}\label{eq:mahlercoef}
 \norm T_2\le\sqrt{\binom{2\ell}{\ell}}\,\mathcal M(T)
            \le2^\ell\mathcal M(T),\qquad \ell=\deg T.
\end{equation}
Indeed each coefficient is bounded by the corresponding binomial
coefficient times \(\mathcal M(T)\), by expanding the product;
sum their squares. See also \cite[Lemma 1.6 and Exercise 1.11]{brunault}.

\begin{lemma}\label{lem:smallproduct}
If \(P\) has degree \(n\) and \(H\ne0\) has degree at most \(d\), then
\[
 \norm{HP}_2\ge\norm H_2(n+d+2)^{-d}\norm P_2.
\]
\end{lemma}
\begin{proof}
For \(V=\sum_{j=0}^Na_jz^j\) and \(\gamma\in\C\), Parseval gives
\[
 \norm{(z-\gamma)V}_2^2=(1+|\gamma|^2)\norm V_2^2
 -2\Rea\left(\overline\gamma
            \sum_{j=0}^{N-1}a_j\overline{a_{j+1}}\right).
\]
The numerical radius of the truncated shift of dimension \(N+1\)
is \(c_N=\cos(\pi/(N+2))\); this is the classical extremal shift
calculation in \cite{haagerup}. Directly, after a diagonal phase
rotation its real part is the path matrix with off-diagonal entries
\(1/2\). Its eigenvectors have coordinates
\(\sin((j+1)k\pi/(N+2))\), and its largest eigenvalue is \(c_N\).
Thus the squared norm ratio is at least \(1+r^2-2rc_N\), where
\(r=|\gamma|\). For \(r\le1\), completing a square bounds this
below by \(1-c_N^2\); for \(r\ge1\), divide by \(r^2\) and use
the same argument at \(1/r\). Consequently
\[
 \norm{(z-\gamma)V}_2\ge
 \max(1,|\gamma|)\sin\frac{\pi}{N+2}\norm V_2.
\]
Factor \(H\) into \(\ell\le d\) linear factors and apply this
estimate successively at degrees \(n,n+1,\ldots,n+\ell-1\).
Since \(\sin(\pi/x)\ge2/x\) for \(x\ge2\),
\[
 \norm{HP}_2\ge\mathcal M(H)
       \left(\frac2{n+d+2}\right)^\ell\norm P_2
 \ge\norm H_2(n+d+2)^{-\ell}\norm P_2.
\]
This implies the assertion; the constant-polynomial case is immediate.
\end{proof}

\subsection{A positive semicircle in the reciprocal identity}
\begin{proof}[Proof of Theorem~\ref{thm:count}]
Write \(P=cAQ\), where
\(A(z)=\prod_{j=1}^M(z-\beta_j)\), \(|\beta_j|<1\), is monic
and all roots of \(Q\) lie on the circle. For a polynomial considered
with degree parameter \(k\), set
\(T^*(z)=z^k\overline{T(1/\bar z)}\).
Use parameters \(n\) for \(P\), \(M\) for \(A\), and importantly
\(n\) for \(R=P'-\lambda P\), even if \(R\) has smaller actual degree.
Self-inversivity of \(Q\) gives
\(P/P^*=\omega A/A^*\) for some \(|\omega|=1\).
The identity
\(z^2(P^*)'=nzP^*-R^*-\bar\lambda P^*\), followed by
differentiation of this quotient and clearing denominators, yields
\begin{align}
 z^2AA^*R+\omega A^2R^*&=HP,\label{eq:reciprocalH}\\
 H&=z^2\{A'A^*-A(A^*)'\}
       +(nz-\lambda z^2-\bar\lambda)AA^* .\label{eq:Hdef}
\end{align}
These are polynomial identities, including at zeros of the factors.
Here \(\deg H\le d=2M+2\), and on the circle
\(|A^*|=|A|\le2^M\), \(|R^*|=|R|\). Hence
\begin{equation}\label{eq:HPupper}
 \norm{HP}_2\le2\,4^M\norm R_2.
\end{equation}

We obtain a lower bound for \(H\) without a separation assumption.
The finite Blaschke identity, also given in
\cite[Lemma 4.3]{kovalev}, is
\[
 z\frac{(A/A^*)'}{A/A^*}
       =\sum_{j=1}^M\frac{1-|\beta_j|^2}{|z-\beta_j|^2},
       \qquad |z|=1.
\]
It follows either by differentiating each factor or by its Poisson
kernel formula. Since \(AA^*=z^M|A|^2\) on the circle,
\begin{align}
 z^{-M-1}H(z)&=(n-2\Rea(\lambda z))|A(z)|^2+K_A(z),\label{eq:positiveH}\\
 K_A(z)&=\sum_{j=1}^M(1-|\beta_j|^2)
                    \prod_{\ell\ne j}|z-\beta_\ell|^2\ge0.\nonumber
\end{align}
Choose a semicircle \(E\) on which \(\Rea(\lambda z)\le0\);
if \(\lambda=0\), choose either semicircle. Thus \(|H|\ge n|A|^2\)
on \(E\).

Jensen's formula for the monic stable polynomial \(A\) gives
\(\int_{\mathbb T}\log|A|\,dt/(2\pi)=0\).
On the complementary semicircle, \(\log|A|\le M\log2\).
Therefore its average over \(E\), with measure \(dt/\pi\), is
at least \(-M\log2\). Applying Jensen's inequality to
\(\exp(4\log|A|)\) on \(E\) gives
\[
 \frac1\pi\int_E|A|^4\,dt\ge2^{-4M},\qquad
 \norm H_2\ge\frac{n}{\sqrt2\,4^M}.
\]
This restricted-arc Jensen argument is classical in origin; compare
\cite[Lemma 4.1]{defant}. The logarithm is continuous on the circle
for each strictly interior configuration, so no limiting interchange
as the roots approach the boundary is needed.
In particular \(H\ne0\). Lemma~\ref{lem:smallproduct} and
\eqref{eq:HPupper} now give \eqref{eq:count}.
For \(n\ge2M+4\), use \(n+2M+4\le2n\) and weaken
\(2^{-6M-7/2}\) to \(2^{-6M-4}\).
\end{proof}

\begin{corollary}[Necessary quantitative count]\label{cor:complexity}
Under the hypotheses of Theorem~\ref{thm:count},
\[
 M\ge\frac{\log(1/\rho_\lambda(P))-\log n-\log(98\sqrt2)}
              {2\log n+\log784}.
\]
In particular, a sequence with
\(\rho_{\lambda_n}(P_n)\le\exp(-c n^\alpha)\), where
\(c,\alpha>0\), necessarily has
\(M_n\ge(c/2+o(1))n^\alpha/\log n\).
\end{corollary}
\begin{proof}
Since \(M\le n\), we have \(n+2M+4\le7n\).
Taking logarithms in \eqref{eq:count} gives
\[
 \log(1/\rho_\lambda(P))
 \le(2M+1)\log n+M\log784+\log(98\sqrt2).
\]
Rearrangement proves the result. This is only a necessary condition;
Theorem~\ref{thm:main} can also exclude a proposed decay rate.
\end{proof}

\begin{remark}[Stable factor and self-inversive cofactor]\label{rem:structural}
The same bound \eqref{eq:count} holds for any degree-\(n\)
\(P=cAQ\) with \(A\) monic of degree \(M\), all its zeros
strictly inside the disk, and \(Q=\eta Q^*\), \(|\eta|=1\).
Indeed only this last identity was used for \(Q\); the remaining
estimates use the degree of \(P\) and the zeros of \(A\).
The cofactor may therefore contain conjugate-reciprocal pairs
inside and outside the circle. In this extension \(M\) means the
degree of the chosen stable factor, not the total number of interior
zeros of \(P\). Not every polynomial has such a factorization.
For instance, \(P_R=1-z/R\), \(R>1\), with \(\lambda=-1/R\)
has residual \(1/[R^2\sqrt{1+R^{-2}}]\to0\) at fixed degree one;
its unpaired exterior root falls outside the hypothesis.
\end{remark}

\subsection{Sufficiency, necessity and exact saturation}
\begin{proof}[Proof of Theorem~\ref{thm:budget}]
If \(b_n\) does not tend to infinity, there are a finite \(B\)
and infinitely many degrees with \(b_n\le B\).
For every such sufficiently large \(n\), Theorem~\ref{thm:count}
bounds any admissible residual below by
\(2^{-6B-4}n^{-2B-1}\), since its actual count \(M\le B\).
This contradicts superalgebraic decay, so (2) implies (1).

Suppose now \(b_n\to\infty\). For each integer \(j\ge1\),
apply Theorem~\ref{thm:upper} with exponent \(2j\), obtaining
\(M_j,C_j,N_j\). Choose strictly increasing integers \(T_j\ge j\)
so large that for every \(n\ge T_j\),
\[
 b_n\ge M_j,\qquad n\ge N_j,\qquad n^j\ge C_j.
\]
All auxiliary eventual conditions in that construction can also be
included in \(N_j\). With \(j(n)=\max\{j:T_j\le n\}\), the
constructed real normalized polynomial has exactly \(M_{j(n)}\)
interior zeros and restricted residual at most \(n^{-j(n)}\).
The index is finite and tends to infinity.

To achieve exactly \(b_n\) interior zeros while retaining reality,
we do \emph{not} multiply the polynomial by anything. We only move
some of its boundary zeros. Fix \(n\) with \(n\ge T_1\) and put
\(j=j(n)\). Let \(P_n\) be the real polynomial constructed in the proof of
Theorem~\ref{thm:upper} at degree \(n\) with exponent \(2j\); it has degree
exactly \(n\), \(P_n(0)=1\), and
\[
 P_n(z)=B(z)^{\ell}(1+z)^{c}\prod_i f_i(z)^{m_i},
 \qquad B(z)=(1+z)^4(1-z+z^2),
\]
where, in the notation of that proof, \(\ell=\ell_n\) and \(c=c_n\) are the
integer exponents, and \(\prod_if_i^{m_i}\) collects the \(M=M_j\)
interior zeros. The factor \(B^\ell(1+z)^c\) is part of \(P_n\); its
\(6\ell+c=n-M_j\) zeros are the boundary zeros of \(P_n\), namely
\(4\ell+c\) real zeros at \(-1\) and \(\ell\) conjugate pairs
\(\ee^{\pm i\pi/3}\). Eventually \(\ell\ge1\) (all eventual conditions are included
in \(N_j\)). Given the required additional count \(s=b_n-M_j\), so that
\(0\le s\le n-M_j=6\ell+c\), take
\[
 q=\min\{\ell,\lfloor s/2\rfloor\},\qquad r=s-2q .
\]
Then \(0\le q\le\ell\) and \(2q+r=s\). Moreover \(0\le r\le4\ell+c\): if
\(q=\lfloor s/2\rfloor\) then \(r\in\{0,1\}\) and \(4\ell+c\ge4\);
otherwise \(q=\ell\) and \(r=s-2\ell\le4\ell+c\). (We also checked this
bookkeeping by machine for all \(1\le\ell\le24\), \(0\le c\le24\),
\(0\le s\le6\ell+c\).) For \(0<\delta<1\) put \(\varrho=1-\delta\) and
\[
 P_{n,\delta}(z)=(1-z+z^2)^{\ell-q}\bigl(1-z/\varrho+z^2/\varrho^2\bigr)^{q}
 (1+z)^{4\ell+c-r}(1+z/\varrho)^{r}\prod_if_i(z)^{m_i}.
\]
This is \(P_n\) with \(q\) of the \(\ell\) pairs \(\ee^{\pm i\pi/3}\)
replaced by the pairs \(\varrho\ee^{\pm i\pi/3}\) and \(r\) of the zeros
\(-1\) replaced by \(-\varrho\). Its degree is still \(n\), its coefficients are
real, \(P_{n,\delta}(0)=1\), and its zeros lie in the closed disk. For
\(\delta>0\) the \(2q+r=s\) displaced zeros lie strictly inside the disk,
the \(M_j\) interior zeros of \(P_n\) are untouched, and the remaining
zeros stay on the circle; so \(P_{n,\delta}\) has exactly
\(M_j+s=b_n\) interior zeros. As \(\delta\to0\), \(P_{n,\delta}\to P_n\)
coefficientwise. The restricted residual
\(\rho(P)=\norm{P'-P'(0)P}_2/\norm P_2\) is a continuous function of the
coefficients wherever \(\norm P_2\ne0\), and \(\norm{P_n}_2\ne0\). Therefore
\(\rho(P_{n,\delta})\to\rho(P_n)\), and, for each fixed \(n\), we can fix
\(\delta=\delta_n>0\) with \(\rho(P_{n,\delta_n})\le\rho(P_n)+n^{-j(n)}\).
Since \(\rho(P_n)\le C_{j}n^{-2j}\le n^{-j}\) (because \(n^j\ge C_j\)), the
final polynomial has exactly \(b_n\) interior zeros, real coefficients,
\(P(0)=1\), and residual at most \(2n^{-j(n)}=o(n^{-A})\) for every fixed
\(A>0\); the slope used is \(P'(0)\). No product with an extra factor is
taken, and no continuity is needed in \(n\).
The finitely many initial degrees can be filled arbitrarily with
real normalized polynomials having the prescribed count, for example
\((1+z)^{n-b}(1+2z)^{b}\) for \(b\) interior zeros \(-1/2\).
This proves (1) implies (3), while (3) implies (2) immediately.

For residuals merely tending to zero, infinitely many zero budgets
would contradict Theorem~\ref{thm:circle}. Conversely, a budget
at least one eventually accommodates the family of
Theorem~\ref{thm:one}. This last assertion concerns an upper budget,
not saturation of every budget by that particular family.
\end{proof}

For a fixed \(M\), residual \(O(n^{-A})\) consequently requires
\(A\le2M+1\). This necessary relation is not a matching tradeoff.
The number \(M_A\) in Theorem~\ref{thm:upper} depends on \(A\);
no single fixed count gives a superalgebraic sequence.

\subsection{A stronger obstruction for one interior zero}
The reciprocal identity works for every count. At count one, the
third and fourth reciprocal moments give a stronger exponent.

\begin{proposition}\label{prop:one-lower}
If \(n\ge4096\), \(P\in\mathcal S_n\) has exactly one strictly
interior zero and all other zeros on the circle, then, for all
\(\lambda\in\C\),
\[
 \rho_\lambda(P)\ge\frac{n^{-13/6}}{512\sqrt2\,\ee^3}.
\]
\end{proposition}
\begin{proof}
Set \(C=2\sqrt2\,\ee^3\) and suppose for contradiction that
\(\rho=\rho_\lambda(P)\le n^{-13/6}/(256C)\).
Rotate a maximum point to 1 and normalize \(G(1)=\norm G_\infty=1\)
as in Lemma~\ref{lem:normalize}. Its initial part gives
\(|\mu|\le n+1\) and \(S_1=G'(1)\in\mathbb R\), where
\(S_k=\sum_j(1-\beta_j)^{-k}\) uses the rotated zeros.
The local ODE proof of Lemma~\ref{lem:cauchy}, at radius \(1/n\),
does not require its separate separation parameter \(\eta\).
Here its smallness parameter is at most \(\ee^3/n<1/2\), and
it yields the unaveraged estimates
\[
 |S_k|\le2\sqrt{n+1}\rho\ee^{(2n+1)/n}n^{k-1}
       \le C\rho n^{k-1/2},\qquad k\ge2.
\]
In particular,
\begin{equation}\label{eq:one-small}
 |S_2|\le\frac{n^{-2/3}}{256},\quad
 |S_3|\le\frac{n^{1/3}}{256},\quad
 |S_4|\le\frac{n^{4/3}}{256}.
\end{equation}
The \(n-1\) boundary reciprocals are \(w_j=1/2+it_j\),
\(t_j\in\mathbb R\); the remaining reciprocal is \(v=x+iy\),
\(x>1/2\). Since \(S_1\) is real, \(\sum_jt_j=-y\), so
\[
 \operatorname{Im}S_2=(2x-1)y,
\]
and, using \(\Rea w_j^3=(3/2)\Rea w_j^2-1/4\),
\[
 x^3-\tfrac32x^2-3(x-\tfrac12)y^2
   =\frac{n-1}{4}+\Rea S_3-\tfrac32\Rea S_2\ge\frac n8.
\]
The last inequality follows from \eqref{eq:one-small} for
\(n\ge4096\). Both subtracted terms on the left are nonnegative,
so \(x\ge n^{1/3}/2\ge8\). Consequently
\(2x-1\ge x\), \(|xy|\le|S_2|\), and
\[
 \Rea v^4=x^4-6x^2y^2+y^4
 \ge\frac{n^{4/3}}{16}-\frac6{65536}n^{-4/3}.
\]
For each boundary reciprocal,
\(\Rea(1/2+it)^4=(t^2-3/4)^2-1/2\ge-1/2\).
It follows that
\[
 \Rea S_4\ge\frac{n^{4/3}}{16}
              -\frac6{65536}n^{-4/3}-\frac{n-1}{2}
            \ge\frac{n^{4/3}}{64},
\]
using \(n/2\le n^{4/3}/32\). This contradicts
\eqref{eq:one-small}.
\end{proof}

Together with Theorem~\ref{thm:one}, this leaves a gap between
the universal one-root lower exponent \(13/6\) and the achieved
upper exponent \(1/6\). No optimality of either is asserted.

\section{A fixed radial gap forces a polynomial barrier}\label{sec:radial}
The upper construction uses the full closed disk. A fixed strict
root radius gives a complementary obstruction, by the same local
estimate and bounded test as the lower-bound proof.

\begin{proposition}\label{prop:radial}
For \(0\le q<1\), let \(s_n(q)\) denote the infimum defining \(s_n\)
with all zeros restricted to \(|z|\le q\). Set
\[
 a=1-q^2,\quad K=\lceil4/a\rceil,\quad
 \Delta=\frac1{1+q}-\frac{1+q^K}{Ka}>0,\qquad
 c_q=\min\left\{\frac1{\sqrt2},
       \frac{\Delta Ka}{2\sqrt2\ee^3(1+a)^K}\right\}.
\]
Then \(s_n(q)\ge c_q n^{3/2-K}\) for every \(n\ge41\).
\end{proposition}
\begin{proof}
Write \(\rho\) for an admissible residual. If
\(\rho>1/\sqrt{n+1}\), the conclusion is immediate.
Otherwise the maximum normalization gives \(|\mu|\le n+1\).
Take \(r=1/n\). The smallness parameter in the proof of
Lemma~\ref{lem:cauchy} is at most \(\ee^3/n<1/2\).
That local proof does not need the separate zero-separation
parameter \(\eta\), and gives
\[
 |m_j|\le\frac{2\sqrt{n+1}}n\rho\ee^3 r^{1-j}\qquad(j\ge2).
\]
The exact image of \(|\beta|\le q\) under \(v=(1-\beta)^{-1}\)
is \(|v-1/a|\le q/a\). Thus \(|1-av|\le q\) and
\(\Rea v\ge1/(1+q)\).
Use \(\Psi(v)=[1-(1-av)^K]/(Ka)\), whose modulus is at most
\((1+q^K)/(Ka)\). Estimating its higher averaged terms gives
\[
 \Delta\le\frac{2\ee^3\sqrt{n+1}\rho}{Kan^2}
       [(1+an)^K-1-Kan]
 \le\frac{2\ee^3\sqrt{n+1}\rho}{Kan^2}(1+an)^K.
\]
Now \(\sqrt{n+1}\le\sqrt{2n}\) and \(1+an\le(1+a)n\).
The resulting lower bound is the stated one. Finally \(Ka\ge4\),
\(q^K<1\), and \(1/(1+q)>1/2\) ensure \(\Delta>0\).
\end{proof}

In particular, any single sequence of admissible degree-\(n\)
polynomials with residual \(o(n^{-A})\) for every fixed \(A>0\)
must have its maximum root modulus tending to one. Otherwise an
infinite subsequence would satisfy a fixed polynomial lower bound.
No sharp radial exponent is asserted.

\section{Scope and verification}
The uniform slope formulation bounds approximate eigenvectors of
polynomial differentiation under a zero-location constraint.
It does not describe the full unconstrained pseudospectrum.
For example, for \(u=n/4\) and
\(T_n(z)=\sum_{k=0}^n(uz)^k/k!\),
\[
 T_n'-uT_n=-\frac{u^{n+1}}{n!}z^n,\qquad
 \frac{\norm{T_n'-uT_n}_2}{\norm{T_n}_2}
       \le\frac n4\left(\frac{\ee}{4}\right)^n,
\]
using \(\norm{T_n}_2\ge1\) and \(n!\ge(n/\ee)^n\).
Thus a root constraint is essential for a bound of our form.

The lower and upper disk bounds do not match in order.
Whether the restricted exterior power-sum minima stay bounded
away from zero remains unresolved here, as does the optimal order
of \(s_n\) and \(r_n\).
The superalgebraic upper bound is not a stretched-exponential upper
bound, and its constants are not optimized. The minimal interior-root
count permitting collapse is one, but the optimal rate with one such
root is not established. Theorem~\ref{thm:count} excludes
superalgebraic decay for every fixed number of interior roots. The elementary radial
barrier is a consequence of the same local method.
We do not assert a theorem for arbitrary starting exponents,
weighted power sums, or \(H^p\) with \(p<1\).
The elementary norm comparison in Corollary~\ref{cor:hp} is a
consequence of the main estimate.

The analytic proof of the explicit threshold uses only the displayed
elementary inequalities. It does not rely on optimization or sampling.
The supplementary diagnostics check exact projected Newton identities
over Gaussian rationals, the bounded test on admissible reciprocal
zeros, binomial identities, and numerical instances of the finite
certificate. Rational checks of constants and high-precision
evaluations are included for reproducibility. The circle-root family
is checked with complete exact coefficient Parseval ratios.
The one-interior-root family has symbolic and high-precision checks
of the cubic, slope, moments and circle parameters.
A separate calculation by an AI model that also took part in preparing
this manuscript, hence not independent, covers
all 34 degree classes at two scales.
Additional exact checks cover the reciprocal identity, its complex
conjugations, the positive Poisson kernel, the one-root moment
obstruction and real-root saturation.

For Proposition~\ref{prop:prefactor}, supplementary numerical norm
diagnostics use two equispaced grids, of sizes \(n+1\) and \(n+2\),
at eight degrees up to \(10\,000\,017\). Such grids would integrate
the squared polynomial norms exactly in exact arithmetic. The
implemented evaluations instead use rounded parameters and omit
underflow-scale weights. The cancellation-free remainder approximates
the exact family; it is not literally the exact residual of the
polynomial with all parameters rounded. Agreement between grids
does not certify the common rounding or masking error.
These computations are not interval-certified norm evaluations and
do not prove the asymptotic limit. Finite moment searches motivated
the general construction but are not premises of its arbitrary-order
proof. None of the checks establishes historical priority or replaces
the analytic arguments.

In version 2, the proofs of Theorem~\ref{thm:upper}, Theorem~\ref{thm:one}
and Proposition~\ref{prop:prefactor} were rewritten with the uniform
estimates displayed (Lemmas~\ref{lem:cone}, \ref{lem:factors} and~\ref{lem:assemble}). As a
consistency check, not a premise of any proof, one explicit instance of the
construction of Theorem~\ref{thm:upper} with \(K=3\) and
\(M=10\) interior zeros (three real groups and one conjugate pair, with small
integer multiplicities satisfying the moment conditions of
Lemma~\ref{lem:cone}) was solved, with 60-digit parameters, for
\(n=10^5,10^6,10^7\): the identities \(S_1=P_n'(0)\) and
\(S_2=S_3=0\) held with absolute errors below \(10^{-33}\), and the residual, computed in double
precision on equispaced grids of sizes about \(n\) (two grid sizes agreeing to
about ten digits), satisfied \(n^{1/6}\rho(P_n)\approx282,244,238\), that is,
decay like \(n^{-1/6}\) with a large constant, in agreement
with the bound \(n^{(2-K)/6}\). No instance with
\(K\ge4\) was computed, because the integer weights produced by
Lemma~\ref{lem:cone} for \(K=4\) were too large for the same tests. For
Theorem~\ref{thm:one} the same tests, up to \(n\approx10^7\), gave
\(n^{1/6}\rho(P_n)\) decreasing slowly towards \(0.61\), consistent with the
constant of Proposition~\ref{prop:prefactor}; this is again not a proof.

The comparison uses the exact primary-source versions listed below.
A targeted search did not identify the bounds in this exact setting;
exhaustive priority certification is not claimed. The accessible
Komarov versions listed below were used for the stated comparison.
A related 2020 journal text (DOI~\href{https://doi.org/10.1070/IM8901}{10.1070/IM8901})
was not available in full in this review;
no exclusion of every possible bibliographic overlap is claimed.
AI tools assisted exploration, drafting and computational checking.
Model review is distinct from independent human peer review and
formal proof verification.

\begin{thebibliography}{99}
\bibitem{tan}
X. Tan, Q. Wang, W. Huang and K. Chen,
\emph{Residual bounds for Schur-stable polynomials},
arXiv:2608.02043v3, 8 August 2026.
\url{https://arxiv.org/abs/2608.02043v3}.
\bibitem{luo}
Y. Luo, R. Yang and K. Zhu,
\emph{Exterior power sums},
arXiv:2607.22017v1, 24 July 2026.
\url{https://arxiv.org/abs/2607.22017v1}.
\bibitem{komarovrate}
M. A. Komarov,
\emph{Rate of approximation by logarithmic derivatives of polynomials
whose zeros lie on a circle}, arXiv:1804.08918v1, 24 April 2018.
\url{https://arxiv.org/abs/1804.08918v1}.
\bibitem{komarovmoments}
M. A. Komarov,
\emph{Caratheodory type representation with unit weights and related
approximation problems}, arXiv:1807.06499v1, 17 July 2018 (Russian).
\url{https://arxiv.org/abs/1807.06499v1}.
\bibitem{erdelyi}
T. Erd\'elyi, D. P. Hardin and E. B. Saff,
\emph{Inverse Bernstein inequalities and min-max-min problems on
the unit circle}, arXiv:1307.4056v1, 15 July 2013.
\url{https://arxiv.org/abs/1307.4056v1}.
\bibitem{busch}
P. Busch, J. Kiukas and R. F. Werner,
\emph{Sharp uncertainty relations for number and angle},
arXiv:1604.00566v2, \S VI.A.2.
\url{https://arxiv.org/abs/1604.00566v2}.
\bibitem{dlmf}
NIST Digital Library of Mathematical Functions,
\S\S28.2 and 28.8, especially Eq.~(28.8.1).
\url{https://dlmf.nist.gov/28.2};\newline
\url{https://dlmf.nist.gov/28.8}. Accessed 18 October 2026.
\bibitem{brunault}
F. Brunault and W. Zudilin,
\emph{Many Variations of Mahler Measures: A Lasting Symphony},
Australian Mathematical Society Lecture Series 28,
Cambridge University Press, 2020, Chapter 1.
\url{https://doi.org/10.1017/9781108885553}.
\bibitem{haagerup}
U. Haagerup and P. de la Harpe,
\emph{The numerical radius of a nilpotent operator on a Hilbert space},
Proc. Amer. Math. Soc. \textbf{115} (1992), 371--379.
\url{https://doi.org/10.1090/S0002-9939-1992-1072339-6}.
\bibitem{defant}
A. Defant and M. Masty\l o,
\emph{\(L^p\)-norms and Mahler's measure of polynomials on the
\(n\)-dimensional torus}, Constr. Approx. \textbf{44} (2016), 87--101.
\url{https://doi.org/10.1007/s00365-015-9319-x}.
\bibitem{kovalev}
L. V. Kovalev and X. Yang,
\emph{Extreme values of the derivative of Blaschke products and
hypergeometric polynomials}, arXiv:2007.09760v2, 25 June 2021;
Bull. Sci. Math. \textbf{169} (2021), 102979.
\url{https://arxiv.org/abs/2007.09760v2}.
\end{thebibliography}
\end{document}
